\documentclass{article}
\usepackage[T1]{fontenc}
\usepackage[utf8]{inputenc}
\usepackage{times}
\usepackage[french]{babel}
\usepackage[left=0.7in,right=0.7in,top=0.4in,bottom=1.0in]{geometry}

\usepackage{import}

\input{../LatexUtils/packages.tex}
\input{../LatexUtils/colors.tex}
\input{../LatexUtils/macro.tex}
\input{../LatexUtils/boxes.tex}
\input{../LatexUtils/customlst.tex}
 
\title{Spécification et preuve formelle de programmes : ACSL (TD3)}
\date{}

\begin{document}

\maketitle

\vspace*{-0.5in}

\begin{Memo}{Variant et behaviors}{}
    Frama-C est un outil puissant contenant de nombreuses fonctionnalités.
    Nous n'avons vu ici qu'une petite partie de ses capacités, notamment les 
    modules WP et RTE.
    Et même parmi ces modules, nous n'avons vu qu'une partie des possibilités offertes.
    Il reste notamment de nombreuses annotations ACSL que nous n'avons pas abordées.
    Par exemple, à mon grand dam, la plupart de vos versions Frama-C ajoutent
    automatiquement une clause vis-à-vis de la terminaison des fonctions 
    (\texttt{\textbackslash terminates}).
    Et c'est en effet un aspect important de la correction d'un programme.
    Ainsi, pour prouver qu'une fonction termine, il est nécessaire, en particulier,
    de prouver que toutes les boucles qu'elle contient terminent.
    Pour ce faire, on utilise des \textbf{variants de boucle}.
    Un variant de boucle est une expression entière qui doit être :
    \begin{itemize}[label=$\bullet$]
        \item Toujours positive au début de chaque itération de la boucle ;
        \item Décroissante à chaque itération de la boucle.
    \end{itemize}
    S'il est possible de prouver l'existence d'un tel variant de boucle,
    alors on peut naturellement conclure que la boucle termine,
    puisque le variant possède un nombre fini de valeurs possibles.
    Ainsi, dans l'exemple suivant, on peut utiliser le variant de boucle
    \texttt{n - i} pour prouver que la boucle termine.
    En ajoutant l'annotation \texttt{\textbackslash terminates} à la fonction,
    on peut tester avec Frama-C, et vérifier que la preuve de terminaison
    passe correctement.
    \begin{MeCode}
        \begin{lstlisting}[style=CStyle]
/*@ terminates \true; */
int sum(unsigned int n) {
    unsigned int res = 0;
    unsigned int i = 0;
    /*@ loop variant n - i */
    while(i < n) {
        res += i;
        i++;
    }
    return res;
}
        \end{lstlisting}
    \end{MeCode}

    \medskip

    Un autre outil très utile est la possibilité de définir des \textbf{behaviors} (comportements)
    pour une fonction.
    Il arrive souvent qu'en fonction des arguments passés à une fonction,
    le comportement de cette fonction change.
    Par exemple, dans le code ci-dessous, la fonction effectue soit une addition,
    soit une soustraction en fonction de la valeur de \texttt{a}.

    \begin{MeCode}
        \begin{lstlisting}[style=CStyle]
Status perform_action(Action a, int x, int y, int * res) {
    switch(a) {
        case ADD: if((y >= 0 && x > INT_MAX - y) || (y <= 0 && x < INT_MIN - y)) return FAIL; 
                  else { *res = x + y; return SUCCESS; }
        case SUB: if((y >= 0 && x < INT_MIN + y) || (y <= 0 && x > INT_MAX + y)) return FAIL; 
                  else { *res = x - y; return SUCCESS; }
        default: return FAIL;
    }
    return FAIL;
}
        \end{lstlisting}
    \end{MeCode}

    On \textit{pourrait} écrire une spécification ACSL pour cette fonction
    avec les outils vus jusqu'à présent, mais cela serait fastidieux et peu lisible.
    On peut alors utiliser des behaviors afin de décrire le comportement de la fonction
    en fonction de la valeur de ses arguments, par exemple de \texttt{a}.
    Un exemple est donné dans le fichier \texttt{TD3/src/behaviors\_example.c}.
    On remarque que chaque behavior est constituée d'une précondition
    décrivant le cas où ce comportement s'applique, par exemple 
    \texttt{assumes a == ADD;} pour le comportement d'addition.
    De plus chaque behavior possède aussi une postcondition décrivant
    l'effet de la fonction dans ce cas.
    Enfin, notons les deux dernières lignes \texttt{complete behaviors add, sub, fail;} et
    \texttt{disjoint behaviors add, sub, fail;}.
    Si la première est prouvée, cela indique que les comportements 
    listés couvrent tous les cas possibles
    (c'est-à-dire que pour toute valeur des arguments, au moins un des comportements
    s'applique).
    Si la seconde est prouvée, cela indique que les comportements 
    listés sont mutuellement exclusifs
    (c'est-à-dire que pour toute valeur des arguments, au plus un des comportements
    s'applique).
    Lorsqu'on définit des behaviors, il est préférable de vérifier que ces deux 
    propriétés sont vraies, afin d'éviter des comportements non couverts
    ou des comportements qui se chevauchent.

    \medskip

    Il y a une différence fondamentale entre le mot clé \texttt{assumes}
    et le mot clé \texttt{requires} vu précédemment (qui peut aussi
    apparaître dans un behavior).
    En effet, une précondition \texttt{requires P;} indique que la fonction
    ne doit être appelée que si la propriété \texttt{P} est vraie.
    Tandis qu'une précondition \texttt{assumes P;} indique que
    le comportement décrit par le behavior s'applique uniquement
    lorsque la propriété \texttt{P} est vraie.
    Ainsi, si on utilise le mot clé \texttt{assumes}, on exprime le comportement
    de la fonction si la propriété est vraie, mais on ne dit pas 
    implicitement qu'il n'existe pas un autre comportement si la propriété
    est fausse.
    Tandis qu'avec le mot clé \texttt{requires}, on exprime que
    la fonction ne doit pas être appelée si la propriété est fausse, et donc
    qu'il n'existe pas d'autre comportement dans ce cas.
    On peut observer cette différence dans le fichier \texttt{TD3/src/requires\_vs\_assumes.c}.
    En analysant ce fichier avec Frama-C, on remarque que la preuve
    de la fonction \texttt{div\_2} échoue, cela vient du fait que du point de vue de
    Frama-C, il existe des cas non considérés (par exemple lorsque \texttt{y == 0}).
\end{Memo}

\begin{Memo}{Prédicats et fonctions logiques}{}
    Comme dans tout langage de programmation, il est possible de 
    définir des sortes de fonctions ou variables, qui permettent
    d'éviter la redondance.
    Par exemple, il est possible de définir des \textbf{prédicats}
    en ACSL, qui sont des formules logiques nommées, que l'on peut
    réutiliser dans les annotations ACSL.
    Par exemple, le prédicat suivant prend en entrée deux entiers \texttt{x} et \texttt{y}, 
    et renvoie \texttt{true} si la division de \texttt{x} par \texttt{y} ne cause
    pas d'overflow (ou d'underflow) :
    \begin{MeCode}
        \begin{lstlisting}[style=CStyle]
/*@ predicate valid_div(integer a, integer b) =
        (a != INT_MIN || b != -1) && b != 0;
*/

/*@ requires valid_div(x, y);
    ensures \result == x / y;
*/
int div(int x, int y) {
    return x / y;
}
        \end{lstlisting}
    \end{MeCode}

    Cet exemple peut être testé dans le fichier \texttt{TD3/src/predicate\_example.c}.
    Il est aussi possible de donner des labels aux prédicats, afin d'exprimer
    la temporalité.
    Par exemple, le prédicat ci-dessous permet de comparer la valeur de $*x$
    à moment $L$ et à moment $M$.

    \begin{MeCode}
        \begin{lstlisting}[style=CStyle]
/*@ predicate cmp{L, M}(int * x) =
        \at(*x, L)  == \at(*x, M);
*/\end{lstlisting}
    \end{MeCode}

    Quelques valeurs possibles de $L$ et $M$ sont\footnote{Il en existe quelques autres,
    mais dans le cadre de ce cours, elles ne sont pas nécessaires.} :
    \begin{itemize}[label=$\bullet$]
        \item \texttt{Pre} : avant l'exécution de la fonction ;
        \item \texttt{Here} : à l'endroit courant dans le code ;
        \item \texttt{Post} : après l'exécution de la fonction.
    \end{itemize}
    Ainsi, si on utilise le prédicat \texttt{cmp\{Pre, Post\}(x)} dans la spécification
    d'une fonction, on exprime que la valeur pointée par \texttt{x}
    ne change pas durant l'exécution de la fonction.
    Un exemple est donné dans le fichier \texttt{TD3/src/temporal\_predicate\_example.c}.

    \medskip

    De la même manière, il est possible de définir des \textbf{fonctions logiques}
    en ACSL, qui sont des fonctions nommées que l'on peut réutiliser
    dans les annotations ACSL.
    Par exemple, la fonction logique suivante calcule la valeur absolue
    d'un entier :
    \begin{MeCode}
        \begin{lstlisting}[style=CStyle]
/*@ logic integer abs(integer x) =
        (x >= 0 ? x : -x);
*/

/*@ requires x != INT_MIN;
    ensures \result == abs(x);
*/
int abs(int x) {
    if(x >= 0) return x;
    return -x;
}
        \end{lstlisting}
    \end{MeCode}

    Cet exemple est disponible dans le fichier \texttt{TD3/src/logic\_example.c}.
\end{Memo}

\begin{Exercice}{}{}
    Compléter le fichier \texttt{TD3/src/ex1.c} afin que Frama-C puisse prouver
    la spécification que vous aurez écrite.

    \medskip

    Une fois que Frama-C est en mesure de prouver votre spécification, 
    utiliser exactement la même spécification dans le fichier \texttt{ex1\_bonus.c},
    et assurez-vous de ne rien avoir oublié.
\end{Exercice}

\begin{Exercice}{}{}
    Donner une spécification à la fonction dans \texttt{TD3/src/ex2.c}.
    Il est probable que dans les invariants de boucle, vous souhaitiez
    parler de la valeur des cases du tableau \texttt{a} avant la boucle.
    Pour ce faire, vous pouvez utiliser le mot clé
    \texttt{\textbackslash at(..., Pre)} (car \texttt{\textbackslash old} 
    n'est pas défini en dehors de la spécification d'une fonction).

    \medskip

    De plus, il est possible que les tableaux se chevauchent (d'un point de vue 
    théorique, en pratique, on est bien d'accord que ce n'est pas une bonne idée de faire ça).
    Cependant, pour Frama-C cela est tout à fait possible et pour faire une preuve complète,
    il est nécessaire de le prendre en compte.
    Pour spécifier que les tableaux ne se chevauchent pas, vous pouvez utiliser
    l'annotation \texttt{\textbackslash separated(...)}.
    Par exemple, pour indiquer que les tableaux \texttt{a} et \texttt{b}
    de taille \texttt{n} et \texttt{m} respectivement
    ne se chevauchent pas, vous pouvez écrire :
    \[
        \texttt{\textbackslash separated(a + (0 .. n - 1), b + (0 .. m - 1))}
    \]
    Enfin, une autre notation qui peut vous être utile est celle permettant de
    désigner une plage d'indices dans un tableau (ce sera utile pour indiquer
    quelles parties du tableau sont modifiées par la boucle) :
    \[
        \texttt{loop assigns a[0 .. n - 1]}
    \]
\end{Exercice}

\begin{Exercice}{}{}
    Compléter le fichier \texttt{TD3/src/ex3.c} afin que Frama-C puisse prouver
    la spécification que vous aurez écrite.
    

    \medskip

    Une fois cela fait, essayer d'utiliser votre même spécification dans le fichier
    \texttt{ex3\_bonus.c}.
    Il y a de fortes chances que tout soit validé, pourtant la fonction
    \texttt{sort} ne trie pas correctement le tableau (il réécrit la première case).
\end{Exercice}

\begin{Exercice}{}{}
    Un problème courant avec les expressions logiques est 
    qu'elles sont difficiles à utiliser dans des preuves pour Frama-C, 
    en particulier lorsqu'elles sont récursives.
    Pour observer ce phénomène, analyser le fichier \texttt{TD3/src/sum\_example.c}
    avec Frama-C.
    Constatez que l'analyse n'est pas concluante (je vous laisse vous convaincre
    à votre rythme qu'elle devrait l'être).

    \medskip

    Le problème vient du fait que les expressions logiques récursives
    sont difficiles à manipuler pour les prouveurs automatiques.
    En effet, pour utiliser une telle expression, le prouveur doit
    la dérouler totalement (ce qui est en général infaisable pour des entrées
    de taille non bornée).
    Pour résoudre ce problème, on peut utiliser des \textbf{axiomes} pour définir
    le comportement d'une expression logique sans expliquer
    comment la calculer.
    Cela permet aux prouveurs automatiques d'utiliser
    directement les propriétés de l'expression logique, sans avoir
    à la dérouler.
    Dans le même fichier, essayez d'enlever l'espace avant l'arobase
    à la ligne \texttt{axiomatic SumTab ...} (et pensez à ajouter un
    espace avant la définition de la fonction logique \texttt{partial\_sum}
    qui est juste au-dessus).
    En relançant l'analyse, vous verrez que Frama-C arrive à prouver la
    fonction \texttt{sum} correctement.

    \medskip

    Les axiomes sont un outil puissant mais :
    \begin{itemize}[label=$\bullet$]
        \item Il n'est pas facile de savoir quels axiomes sont utiles ;
        \item Les axiomes sont des propriétés non prouvées (comme un vrai
        axiome en mathématique), ainsi si vous écrivez n'importe quoi,
        Frama-C pourra le réutiliser, ce qui peut mener à des preuves
        incorrectes.
    \end{itemize}
    Par exemple, l'axiome \texttt{SumRec} exprime le fait que la somme
    des valeurs entre les indices \texttt{start} et \texttt{end} est égale
    à la somme des valeurs entre \texttt{start} et \texttt{end - 1},
    plus la valeur à l'indice \texttt{end - 1}.
    Concrètement, si \texttt{t = [1, 2, 3, 4, 5]}, \texttt{start = 1} et 
    \texttt{end = 3}, ce que dit l'axiome, c'est que la somme
    \texttt{t[1] + t[2]} est égale à \texttt{t[1] + t[2] + t[3] - t[3]}.
    Clairement, cette propriété est vraie, mais on aurait pu écrire :
    \[
        partial\_sum(t, start, end) ==  partial\_sum(t, start, end + 1)
    \]
    et là clairement, ce n'est pas vrai (sauf si \texttt{t[end] = 0}).
    Pourtant, Frama-C aurait pu utiliser cet axiome incorrect
    pour prouver des propriétés fausses\footnote{En pratique,
    les gens sérieux (= pas nous dans ce cours) n'écrivent pas des axiomes
    gratuitement.
    Ils les prouvent avec des outils adaptés, par exemple en Coq.}.

    \medskip

    La vraie question est alors : comment savoir quels axiomes sont
    utiles (et idéalement corrects) ?
    En général, c'est assez difficile, mais une intuition que vous pouvez 
    avoir est que les axiomes sont utiles pour représenter des propriétés
    récursives.
    Et ces dernières sont utiles dans les boucles.
    Par exemple, toujours dans le même fichier, on a l'invariant de boucle :
    \[
        \texttt{result == partial\_sum(t, 0, i)}
    \]
    Si on calcule la WLP de la fin de la boucle vers le début, avec cet invariant,
    on obtient :
    \[
        \texttt{result + t[i] == partial\_sum(t, 0, i + 1)}
    \]
    Ainsi, on voit que Frama-C aura besoin de vérifier cette propriété,
    mais la partie droite contient notre expression logique récursive,
    donc Frama-C aura du mal à la prouver sans axiome.
    Ainsi, un axiome intéressant ici, serait celui qui exprime exactement
    cette propriété (c'est \texttt{SumRec} en l'occurrence).

    \medskip

    Appliquer les axiomes au fichier \texttt{TD3/src/ex4.c}.
\end{Exercice}

\begin{Exercice}{}{}
    Compléter le fichier \texttt{TD3/src/ex5.c}.
\end{Exercice}

% \begin{Exercice}{}{}
%     Modifier le fichier \texttt{example.c} pour ajouter la bonne précondition 
%     à la fonction \texttt{div}, et tester à nouveau la preuve avec Frama-C (pensez à bien
%     insérer les gardes avec \texttt{Insert wp-rte guards} avant de prouver).
%     Faites le même test pour la fonction \texttt{add\_one}, et modifier la 
%     précondition si nécessaire.

%     \medskip

%     Pour gagner du temps, vous pouvez lancer Frama-C en lui demandant en avance d'ajouter
%     les gardes et de faire les preuves de toutes les annotations.
%     Pour ce faire, utilisez la commande \texttt{frama-c-gui -wp -wp-rte -wp-prover z3,alt-ergo <fichier>}.
% \end{Exercice}

% \begin{CorrExercice}{}{}
%     Voir le fichier \texttt{TD2/src/example\_correction.c}.
%     On remarquera que la fonction \texttt{add\_one} n'avait pas besoin
%     de modification.
% \end{CorrExercice}

% \begin{Exercice}{}{}
%     Écrire la spécification ACSL des fonctions \texttt{mod} et \texttt{mod2}
%     du fichier \texttt{TD1/src/ex1.c} et vérifier la correction de ces fonctions
%     avec Frama-C.
% \end{Exercice}

% \begin{CorrExercice}{}{}
%     Voir le fichier \texttt{TD2/src/ex1\_correction.c}.
% \end{CorrExercice}

% \begin{Exercice}{}{}
%     Écrire la spécification ACSL de la fonction \texttt{max}
%     du fichier \texttt{TD1/src/ex2.c} et vérifier la correction de cette fonction
%     avec Frama-C.
% \end{Exercice}

% \begin{CorrExercice}{}{}
%     Voir le fichier \texttt{TD2/src/ex2\_correction.c}.
% \end{CorrExercice}

% \begin{Exercice}{}{}
%     Écrire la spécification ACSL de la fonction \texttt{div}
%     du fichier \texttt{TD1/src/ex4.c} et vérifier la correction de cette fonction
%     avec Frama-C.
% \end{Exercice}

% \begin{CorrExercice}{}{}
%     Voir le fichier \texttt{TD2/src/ex4\_correction.c}.
% \end{CorrExercice}

% \begin{Exercice}{}{}
%     Écrire la spécification ACSL de la fonction \texttt{swap}
%     du fichier \texttt{TD1/src/ex3.c} et vérifier la correction de cette fonction
%     avec Frama-C.

%     \medskip

%     Vous aurez sûrement envie de préciser que les pointeurs passés en argument
%     sont \textit{valides} (typiquement non nuls\footnote{En réalité, un pointeur
%     peut être invalide sans être nul, par exemple après un free.}).
%     Pour décrire cette propriété, vous pouvez utiliser l'annotation
%     \texttt{\textbackslash valid(...)} qui permet de garantir qu'un pointeur est valide.
% \end{Exercice}

% \begin{CorrExercice}{}{}
%     Voir le fichier \texttt{TD2/src/ex3\_correction.c}.
%     On remarque qu'on a besoin de deux nouvelles annotations ACSL :
%     \begin{itemize}[label=$\bullet$]
%         \item Le mot clé \texttt{\textbackslash old} permet de faire référence à la valeur d'une variable
%         avant l'exécution de la fonction.
%         \item Le mot clé \texttt{\textbackslash valid} permet de garantir qu'un pointeur est valide (qu'il pointe vers une adresse mémoire allouée).
%     \end{itemize}
% \end{CorrExercice}

% \begin{Exercice}{}{}
%     Ouvrez le fichier \texttt{TD2/src/ex5.c}, dedans vous trouverez la fonction 
%     \texttt{sum}.
%     Cette fonction calcule la somme des entiers de $0$ à $n$.
%     Dans le fichier, la spécification est déjà donnée.

%     \medskip

%     \textbf{1.} Tenter de prouver la spécification avec Frama-C.
%     Que constatez-vous ?

%     \medskip

%     En fait, la méthode que nous avons vue pour calculer la weakest precondition
%     ne fonctionne pas telle quelle pour les boucles (et c'est bien normal, nous n'avons
%     pas vu comment traiter le cas des boucles dans nos règles de calcul de weakest precondition).
%     Dans les faits calculer la weakest precondition d'une boucle est pénible.
%     Ce qui pose problème, c'est que pour calculer l'effet d'une boucle, il faudrait
%     savoir combien de fois la boucle va s'exécuter, or ce n'est pas forcément
%     évident à déterminer statiquement (c'est même en général impossible).
%     Pour résoudre ce problème, on aimerait bien avoir une sorte de "résumé" de la boucle,
%     qui nous dit, sans avoir à dérouler toute la boucle, quel est son effet.
%     Par exemple, dans la fonction \texttt{sum}, on aimerait bien dire :
%     "\textit{après la boucle, la variable \texttt{res} vaut la somme des entiers de $0$ à $n$}".

%     \medskip

%     Pour ce faire, on va utiliser des \textbf{invariants de boucle}.
%     Un invariant de boucle est une propriété qui doit être vraie :
%     \begin{itemize}[label=$\bullet$]
%         \item Au début de chaque itération de la boucle ;
%         \item A la fin de chaque itération de la boucle ;
%         \item Après la fin de la boucle.
%     \end{itemize}
%     Ainsi, lors du calcul de la weakest precondition, on pourra utiliser
%     cet invariant pour résumer l'effet de la boucle, puisque cet invariant
%     est en particulier vrai au début de la boucle (ce qui permet de calculer la weakest precondition
%     avant la boucle).
%     La partie compliquée est de trouver le bon invariant de boucle, qui permet de prouver
%     la spécification.
%     En effet, un invariant de boucle possible pour la fonction \texttt{sum} est
%     $0 \leq i \leq n + 1$.
%     C'est un invariant tout à fait correct, mais il n'est pas assez fort pour prouver
%     la spécification.
%     On a besoin d'un invariant qui permet d'impliquer la postcondition à la fin de la boucle.

%     \medskip

%     \textbf{2.} Afin d'ajouter un invariant de boucle dans Frama-C, on utilise
%     l'annotation \texttt{loop invariant <formule>;} juste avant la boucle.
%     Un exemple est donné dans le fichier \texttt{ex5\_weak\_invariant.c}.
%     Notons qu'en plus de l'invariant, il est nécessaire d'indiquer
%     les variables modifiées par la boucle avec l'annotation
%     \texttt{loop assigns <liste\_de\_variables>;}.

%     \medskip

%     Essayer de trouver des invariants suffisamment forts pour prouver
%     la spécification de la fonction \texttt{sum} avec Frama-C.
%     Rappelez vous qu'un invariant doit être vrai au début et à la fin
%     de chaque itération de la boucle, ainsi qu'après la fin de la boucle.
% \end{Exercice}

% \begin{CorrExercice}{}{}
%     Voir le fichier \texttt{TD2/src/ex5\_correction.c}.
% \end{CorrExercice}

% \begin{Exercice}{}{}
%     \textbf{1.} Écrire la spécification ACSL de la fonction \texttt{is\_sorted}
%     du fichier \texttt{TD1/src/ex6.c}\footnote{Un tableau est une séquence de pointeurs,
%     donc l'utilisation du mot clé \texttt{\textbackslash valid} est nécessaire pour garantir la validité des pointeurs.}.
    
%     \medskip

%     \textit{Note : Vous aurez sûrement besoin d'un quantificateur, en ACSL, on peut 
%     écrire $\backslash forall ~ integer ~ i; 0 <= i < n ==> i == size;$ pour exprimer
%     "pour tout $i$ entre $0$ et $n - 1$, $i$ est égal à $size$".}

%     \medskip

%     \textbf{2.} Proposer des invariants de boucle permettant de prouver
%     la correction de cette fonction avec Frama-C.
% \end{Exercice}

% \begin{CorrExercice}{}{}
%     Voir le fichier \texttt{TD2/src/ex6\_correction.c}.
% \end{CorrExercice}

% \begin{Exercice}{}{}
%     \textbf{1.} Expliquer ce que fait la fonction \texttt{mystery} dans 
%     le fichier \texttt{TD2/src/ex7.c}.

%     \medskip

%     \textbf{2.} Écrire la spécification ACSL de cette fonction.

%     \medskip

%     \textbf{3.} Proposer des invariants de boucle permettant de prouver
%     la correction de cette fonction avec Frama-C.
% \end{Exercice}

% \begin{CorrExercice}{}{}
%     Voir le fichier \texttt{TD2/src/ex7\_correction.c}.
% \end{CorrExercice}

\end{document}
